% !TEX encoding = UTF-8
\documentclass[a4paper,11pt]{article}

\usepackage[utf8]{inputenc}
\usepackage[francais]{babel}
\usepackage[T1]{fontenc}


%math 
\usepackage{amsmath}
%\usepackage{amssymb} // Decommenter si le dernier n'est pas suffisant

%Chimie
\usepackage[version=3]{mhchem}
%cf http://fr.wikibooks.org/wiki/LaTeX/%C3%89crire_des_formules_chimiques

% Allows for temporary adjustment of side margins
% Lorsque les tableaux sont trop grands
\usepackage{chngpage}
 \usepackage{multirow} 
 
% provides filler text
\usepackage{lipsum}
% usage : \lipsum[1]

% just makes the table prettier (see \toprule, \bottomrule, etc. commands below)
\usepackage{booktabs}

%Mise en page
\usepackage{graphicx}
\usepackage{fancyhdr}
\usepackage[top=3cm, bottom=3cm, left=4cm, right=3cm]{geometry}
\usepackage{listings}
\usepackage{xcolor}
\usepackage{placeins}

\definecolor{gris}{gray}{0.5}
\definecolor{code}{gray}{0.95}

% Table des matières cliquable 
\usepackage{hyperref}
\hypersetup{
    colorlinks, % empécher latex de colorer les liens
    citecolor=black,
    filecolor=black,
    linkcolor=black, % couleur des liens dans la table des matières
    urlcolor=blue
}

%Points dans la table des matières
\usepackage{tocloft}
\renewcommand{\cftsecleader}{\cftdotfill{\cftdotsep}} % Ligne de points dans la table des matières


%Page de garde

\newlength{\larg}
\setlength{\larg}{14.5cm}

% Pour corriger alignement, jouer avec la taille de tabular (± 6cm)
\title{
{\rule{\larg}{1mm}}\vspace{7mm}
\begin{tabular}{p{7,6cm} r}
 \multirow{3}{*}{\includegraphics[width=70px]{vyking-logo2.jpg}} & {\Huge {\bf Langage Vyking}} \\
   & \\
   & {\Large Lexing et parsing}
\end{tabular}\\
\vspace{2mm}
{\rule{\larg}{1mm}}
\vspace{2mm} \\
\begin{tabular}{p{9.8cm} r}
   & {\large \bf Compilateurs} \\
   & {\large \bf INFO0085} \\
   & {\large \bf \bsc{Pr. Pierre Geurts}}\\
   & {\large \bf \bsc{Cyril Soldani}}\\
   & {\large  \today}
\end{tabular}\\
\vspace{10cm}
}
\author{\begin{tabular}{p{13.7cm}}
\bsc{Robin Keunen} s093137\\
\bsc{Pierre Vyncke} s091918\\
1\up{ème} Master Ingénieur Civil
\end{tabular}\\
\hline }
\date{}

\usepackage{times}
\newcommand{\be}{\begin{enumerate}}
\newcommand{\ee}{\end{enumerate}}
\newcommand{\bi}{\begin{itemize}}
\newcommand{\ei}{\end{itemize}}

\definecolor{quotationcolour}{HTML}{F0F0F0}
\definecolor{quotationmarkcolour}{HTML}{1F3F81}

% Double-line for start and end of epigraph.
\newcommand{\epiline}{\hrule \vskip -.2em \hrule}
% Massively humongous opening quotation mark.
\newcommand{\hugequote}{%
  \fontsize{42}{48}\selectfont \color{quotationmarkcolour} \textbf{``}
  \vskip -.5em
}

% Beautify quotations.
\newcommand{\epigraph}[2]{%
  \bigskip
  \begin{flushright}
  \colorbox{quotationcolour}{%
    \parbox{.60\textwidth}{%
    \epiline \vskip 1em {\hugequote} \vskip -.5em
    \parindent 2.2em
    #1\begin{flushright}\textsc{#2}\end{flushright}
    \epiline
    }
  }
  \end{flushright}
  \bigskip
}

\begin{document}

% Titre
\maketitle
\thispagestyle{empty}
\newpage

% Mise en page
\pagestyle{fancy}
\lhead{}
\chead{}
\rhead{\itshape \textcolor{gris}{Vyking - lexing et parsing}}
\lfoot{\itshape \textcolor{gris}{INFO-0085 Compilateurs}}
\cfoot{}
\rfoot{\itshape \textcolor{gris}{\thepage}}
\renewcommand{\headrulewidth}{0.4pt}
\renewcommand{\footrulewidth}{0.4pt}

\newpage 

\section{Introduction}

\section{Trois versions de Vyking}
    Pour séparer les difficultés dans la compilation du langage, nous travaillerons sur trois versions successives du langage.
    La première, \emph{basic\_vyking} n'implémente que les fonctions de bases, le noyau du langage.
    La deuxième, \emph{listed\_vyking} rajoute les listes et les opérateurs de liste.
    La troisième et dernière version \emph{full\_vyking} implémentera les closures.
    
\section{Choix des outils}
    Nous avons isolé quelques lexer/parser dans la multitudes des outils disponibles.
    
    \paragraph{ANTLR}
        ANTLR est capable de générer du code Python pour le Parser, il dispose de beaucoup de documentation.
        L'analyse est top-down, ce qui rend la génération d'erreur plus facile.
        Cependant il vise les grammaires LL, nous ne voulions passer trop de temps à modifier la grammaire.
        De plus, en lisant sur le web, il apparait que le code généré cause beaucoup d'appels de fonctions\cite{antlr_ply_pyparsing}.
        Les appels sont bon marché en Java, le langage de base de ANTLR mais les appels sont chers en Python. 
        D'autres outils sont plus efficaces et orientés Python.
        
    \paragraph{Plex}
        Plex est assez populaire mais cet outil ne fait que l'analyse lexicale.
        Nous préférons une solution intégrée pour ne pas devoir apprendre deux outils.
        
    \paragraph{PLY}
        PLY est une implémentation de lex et yacc pour Python.
        PLY est construit en Python pur, il n'est pas construit sur un noyau C.
        Cette approche le rend plus lent que les outils construit en C mais il est plus efficace (en temps et espace)\cite{PLYtalk} que tous les autres outils Python.
        Il existe beaucoup de ressources en ligne et l'outil fournit des outils de diagnostique.
        Le report d'erreur était un des objectifs de développement pour PLY.        
        D'un point de vue pédagogique, nous avons retenu ce choix car il nous permet de nous familiariser avec Lex et Yacc tout en conservant une approche pythonesque.
        
    \paragraph{} Nous avons retenu PLY, si nous avons assez de temps, nous essaierons de développer nos propres outils de lexing et/ou de parsing
        
        
\section{Lexing de l'indentation}
    Dans une syntaxe à la python, les espaces et l'indentation n'est significative qu'en début de ligne. Les transitions sont illustrées à la figure ??? (CF EVERNOTE bol state transition).
    L'état \verb bol  est enclenché quand le lexer lit un retour à la ligne.
    Une fois dans l'état \verb bol , 3 actions sont possibles : 
    \begin{enumerate}
    \item Lecture d'espace ou de tabulation : le token \verb INDENT  est émis, sa valeur est le nombre d'espaces. Une tabulation vaut 4 espaces par défaut.
    \item Lecture d'un retour à la ligne, le lexer reste dans l'état \verb bol .
    \item Lecture d'un autre caractère, retour à l'état précédent.
    \end{enumerate}

\section{Filtre indentation}

\section{Remarques}
    Les 3 étapes lexing, filtre et parsing sont découplées.
    Mais grâce à l'utilisation de générateur/itérateurs pour construire les classes \verb Lexer  et \verb IndentFilter , les performances sont améliorées.
    En effet, au lieu de passer tout le programme dans le \verb Lexer  et de conserver tous les tokens en mémoire, puis de faire un passe complet à travers le filtre, le \verb Lexer génère les tokens un à un et le filtre ajoute les indent et dedent au vol.
    
    10ème règle de greespun
        
    
            
\bibliography{vyking_bib}
\bibliographystyle{plain}

\end{document}	





















